\chapter{A multilocus route from standing variation to evolving \(\G\)}

\section{Additive allelic effects}

Consider \(L\) diploid biallelic loci. At locus \(\ell\), let the count of the focal allele be
\begin{equation}
X_\ell\in\{0,1,2\}.
\end{equation}
Let \(\mathbf a_\ell\in\R^n\) be the additive phenotypic effect per copy. Relative to a reference mean, the additive phenotype is
\begin{equation}
\mathbf z_A=\sum_{\ell=1}^L(X_\ell-2p_{\ell0})\mathbf a_\ell.
\label{eq:additiveGenotypeModel}
\end{equation}

If the allele frequency in population \(i\) is \(p_{i\ell}\), then under Hardy--Weinberg sampling
\begin{equation}
\E[X_\ell]=2p_{i\ell}.
\end{equation}
Therefore the mean additive phenotype is
\begin{align}
\zbar_i
&=\E[\mathbf z_A]\\
&=\sum_\ell(2p_{i\ell}-2p_{\ell0})\mathbf a_\ell\\
&=\boxed{2\sum_{\ell=1}^L(p_{i\ell}-p_{\ell0})\mathbf a_\ell.}
\label{eq:meanFromAlleleFrequencies}
\end{align}

\inwords The population mean is a weighted sum of allelic effect vectors. Changing an allele frequency moves the mean in the direction of that allele's multivariate effect.

\section{Derive \(\G\) under linkage equilibrium}

At a biallelic diploid locus under Hardy--Weinberg equilibrium,
\begin{equation}
\Var(X_\ell)=2p_{i\ell}(1-p_{i\ell}).
\end{equation}
The covariance contributed by that locus is
\begin{align}
\Cov(X_\ell\mathbf a_\ell)
&=\Var(X_\ell)\mathbf a_\ell\mathbf a_\ell^\top\\
&=2p_{i\ell}(1-p_{i\ell})\mathbf a_\ell\mathbf a_\ell^\top.
\end{align}
If loci are in linkage equilibrium, cross-locus covariances vanish, so contributions add:
\begin{equation}
\boxed{
\G_i=2\sum_{\ell=1}^Lp_{i\ell}(1-p_{i\ell})
\mathbf a_\ell\mathbf a_\ell^\top.}
\label{eq:GfromAlleleFrequencies}
\end{equation}

\begin{geombox}
Each outer product \(\mathbf a_\ell\mathbf a_\ell^\top\) is a rank-one covariance ellipsoid pointing along the allelic effect vector. The full \(\G\) is the weighted sum of those rank-one contributions. Alleles near frequency one half receive the largest weight.
\end{geombox}

With linkage disequilibrium, additional cross-locus terms appear. Schematically,
\begin{equation}
\G_i
=2\sum_\ell p_{i\ell}q_{i\ell}\mathbf a_\ell\mathbf a_\ell^\top
+\sum_{\ell\neq m}\Cov(X_\ell,X_m)\mathbf a_\ell\mathbf a_m^\top.
\end{equation}
The linkage-equilibrium model is therefore a transparent baseline rather than a claim that LD is biologically absent.

\section{Weak-selection allele-frequency dynamics}

Let the local directional selection gradient be \(\betav_i\). The marginal directional effect of one focal allele at locus \(\ell\) is
\begin{equation}
s_{i\ell}=\mathbf a_\ell^\top\betav_i.
\end{equation}
Under weak additive selection, the leading allele-frequency change is
\begin{equation}
\dot p_{i\ell}^{\mathrm{sel}}
=p_{i\ell}(1-p_{i\ell})\mathbf a_\ell^\top\betav_i.
\label{eq:alleleSelectionDynamics}
\end{equation}
With symmetric migration from population \(j\) at rate \(m\),
\begin{equation}
\boxed{
\dot p_{i\ell}
=p_{i\ell}(1-p_{i\ell})\mathbf a_\ell^\top\betav_i
+m(p_{j\ell}-p_{i\ell}).}
\label{eq:alleleDynamicsMigration}
\end{equation}

\section{Recover the multivariate response equation, including migration}

Differentiate \cref{eq:meanFromAlleleFrequencies}:
\begin{equation}
\dot{\zbar}_i=2\sum_\ell\dot p_{i\ell}\mathbf a_\ell.
\end{equation}
Substitute \cref{eq:alleleDynamicsMigration}:
\begin{align}
\dot{\zbar}_i
={}&2\sum_\ell p_{i\ell}(1-p_{i\ell})
\mathbf a_\ell\mathbf a_\ell^\top\betav_i\\
&+2m\sum_\ell(p_{j\ell}-p_{i\ell})\mathbf a_\ell.
\end{align}
The first sum is \(\G_i\betav_i\) by \cref{eq:GfromAlleleFrequencies}. The second is
\begin{align}
2m\sum_\ell(p_{j\ell}-p_{i\ell})\mathbf a_\ell
&=m(\zbar_j-\zbar_i),
\end{align}
provided the populations share the same effect-vector representation and reference. Thus
\begin{equation}
\boxed{
\dot{\zbar}_i
=\G_i\betav_i+m(\zbar_j-\zbar_i).}
\label{eq:LandeWithMigration}
\end{equation}

\begin{cautionbox}
When migration is included at the allele-frequency level, it produces an explicit mean-migration term. The simpler identity \(\dot{\zbar}=\G\betav\) is the selection component or the \(m=0\) limit, not the full dynamics.
\end{cautionbox}

\section{How allele-frequency change alters \(\G\)}

Differentiate \cref{eq:GfromAlleleFrequencies}:
\begin{align}
\dot\G_i
&=2\sum_\ell\frac{\dd}{\dd t}
[p_{i\ell}(1-p_{i\ell})]
\mathbf a_\ell\mathbf a_\ell^\top\\
&=2\sum_\ell[(1-p_{i\ell})-p_{i\ell}]\dot p_{i\ell}
\mathbf a_\ell\mathbf a_\ell^\top\\
&=\boxed{
2\sum_\ell(1-2p_{i\ell})\dot p_{i\ell}
\mathbf a_\ell\mathbf a_\ell^\top.}
\label{eq:GdotAlleles}
\end{align}

\inwords An allele rising toward frequency one half increases its contribution to additive variance; an allele moving from one half toward fixation decreases it. Because different loci point in different multivariate directions, unequal frequency changes can rotate \(\G\) as well as change its total magnitude.

\begin{biobox}
The mechanism does not require new mutation or drift. Divergent selection can reweight pre-existing effect vectors in different populations, causing their covariance ellipsoids and response channels to rotate apart using standing variation alone.
\end{biobox}

\begin{figure}[H]
\centering
\includegraphics[width=0.84\textwidth]{figures/05_standing_variation_rotation.pdf}
\caption{A deterministic multilocus simulation in which divergent selection changes standing allele frequencies, rotates \(\G\), and generates channel misalignment.}
\label{fig:standingVariationRotation}
\end{figure}
